\uuid{iq4j}
\exo7id{2511}
\titre{exo7 2511}
\auteur{mayer}
\organisation{exo7}
\datecreate{2009-04-01}
\isIndication{false}
\isCorrection{false}
\chapitre{Différentiabilité, calcul de différentielles}
\sousChapitre{Différentiabilité, calcul de différentielles}
\cpge{mpsi-mp2i.fonctions-deux-variables.a, pcsi.fonctions-deux-variables.a, ptsi.fonctions-deux-variables.a, mpsi-mp2i.fonctions-deux-variables.b, pcsi.fonctions-deux-variables.b, ptsi.fonctions-deux-variables.b, bcpst1.fonctions-deux-variables.b, bcpst1.fonctions-deux-variables.c, mp-mpi.calcul-differentiel-optimisation.b, pc.fonctions-plusieurs-variables.c, psi.fonctions-plusieurs-variables.c}
\module{Calcul différentiel}
\niveau{L3}
\difficulte{}

\contenu{
\question{
Soit $f:\Rr ^2 \to \Rr$
d\'efinie par
$$f(x,y)=\frac{xy^2}{x^2+y^2} \quad \text{si} \quad  (x,y)\neq (0,0)$$
 et $f(0,0)=0$. Montrer que $f$ est continue sur $\Rr^2$,
que pour tout $u\in \Rr^2 \setminus \{0\}$ $\frac{\partial
f}{\partial u}(0,0)$ existe, mais que $f$ n'est pas
diff\'erentiable en $(0,0)$.
}
}
